% ============================================================
% Justicia Cercana — Diseño del prototipo (CS2H01, Pregunta 1)
% ============================================================
% ANTES DE COMPILAR:
%  1. Completa \nombreEstudiante y \codigoEstudiante abajo.
%  2. Coloca en img/: logo.png y las capturas listadas junto a
%     cada \includegraphics (los nombres ya están definidos).
%  3. Compila con pdflatex (dos veces si algo se ve desalineado):
%     pdflatex main.tex
% ============================================================

\documentclass[11pt,a4paper]{article}

\usepackage[utf8]{inputenc}
\usepackage[T1]{fontenc}
\usepackage[spanish,es-noindentfirst]{babel}
\usepackage{lmodern}
\usepackage[margin=2cm]{geometry}
\usepackage{graphicx}
\usepackage{xcolor}
\usepackage{longtable}
\usepackage{booktabs}
\usepackage{array}
\usepackage{enumitem}
\usepackage{fancyhdr}
\usepackage[hidelinks]{hyperref}
\usepackage{csquotes}
\usepackage{placeins}
\MakeOuterQuote{"}

\setlength{\parindent}{0pt}
\setlength{\parskip}{0.55em}
\renewcommand{\arraystretch}{1.2}

% --- Datos a completar ---
\newcommand{\nombreEstudiante}{NOMBRE APELLIDO}
\newcommand{\codigoEstudiante}{Código de estudiante}

% Figura a prueba de archivos faltantes: si img/#1 no existe, muestra un
% aviso en vez de detener la compilación de todo el documento.
\newcommand{\safefig}[2]{%
  \begin{figure}[!htbp]
    \centering
    \IfFileExists{img/#1}{%
      \includegraphics[width=0.85\linewidth,height=0.72\textheight,keepaspectratio]{img/#1}%
    }{%
      \fbox{\parbox[c][5cm][c]{0.85\linewidth}{\centering\itshape Falta \texttt{img/#1}}}%
    }
    \caption{#2}
  \end{figure}%
}

\pagestyle{fancy}
\fancyhf{}
\fancyhead[L]{Justicia Cercana --- Prototipo de interfaz}
\fancyhead[R]{CS2H01 \textperiodcentered\ Human Computer Interaction}
\fancyfoot[C]{\thepage}
\renewcommand{\headrulewidth}{0.4pt}

\title{Justicia Cercana}
\author{\nombreEstudiante}
\date{\today}

\begin{document}

% ---------------- Portada ----------------
\begin{titlepage}
\thispagestyle{empty}
\begin{center}
\IfFileExists{img/logo.png}{\includegraphics[width=3.2cm]{img/logo.png}}{\fbox{\parbox[c][2.5cm][c]{3.2cm}{\centering\itshape Falta \texttt{img/logo.png}}}}\\[1.2cm]

{\large UTEC --- Departamento de Ciencia de la Computación}\\[0.2cm]
{\large CS2H01 --- Human Computer Interaction}\\[0.2cm]
{\large Examen Parcial --- Pregunta 1: Diseño del prototipo}\\[2.2cm]

{\Huge \textbf{Justicia Cercana}}\\[0.5cm]
{\Large Prototipo de plataforma para juzgados de paz rurales}\\[2.5cm]

{\large \nombreEstudiante}\\
{\large \codigoEstudiante}\\[0.5cm]
{\large \today}

\vfill
{\small Prototipo interactivo construido en React + Tailwind CSS.}\\
{\small \url{https://ihc-parcial-gjdy.vercel.app/}}
\end{center}
\end{titlepage}

\begin{center}
\fbox{\parbox{0.85\linewidth}{\centering
\textbf{Prototipo interactivo en línea:}\\
\url{https://ihc-parcial-gjdy.vercel.app/}
}}
\end{center}
\vspace{0.3cm}

% ---------------- Contexto del caso ----------------
\FloatBarrier
\section{Contexto del caso}

\enquote{Justicia Cercana} plantea el diseño de una aplicación para tableta de 10 pulgadas dirigida a
jueces y juezas de paz que atienden comunidades rurales. Su trabajo comprende recibir a las personas,
registrar casos, dejar constancia de acuerdos o actuaciones notariales, y organizar las actividades que
deben realizar después. Hoy gran parte de esa información se lleva en cuadernos, formularios impresos o
archivos digitales sueltos, lo que provoca duplicidad de datos, dificultad para encontrar un registro
anterior y fechas importantes olvidadas.

El diseño de este prototipo se construyó a partir de tres condiciones reales que el propio caso describe,
y que enmarcan cada decisión tomada:

\begin{itemize}[leftmargin=1.3em]
  \item \textbf{Bajo alfabetismo digital.} Los jueces de paz no necesariamente tienen experiencia usando
  aplicaciones complejas; cada tarea debe guiar al usuario, usar términos familiares y evitar pasos
  innecesarios.
  \item \textbf{Conectividad intermitente como condición normal.} Algunos trabajan en zonas rurales con
  conectividad limitada o inestable y pueden pasar varias horas sin acceso a Internet. La aplicación debe
  permitir seguir trabajando sin perder lo registrado, e indicar con claridad cuándo la información está
  guardada localmente y pendiente de sincronización.
  \item \textbf{Dispositivo táctil de 10 pulgadas.} El diseño debe adaptarse al uso táctil en tableta y,
  además, permitir su uso desde una computadora con conexión a Internet.
\end{itemize}

Sobre estas tres condiciones se construyeron los tres módulos que pide el caso: registro de casos
judiciales, registro de actuaciones notariales, y agenda o calendario.

% ---------------- Enfoque general ----------------
\FloatBarrier
\section{Enfoque general del prototipo}

A diferencia de un mockup estático, este prototipo se construyó como una aplicación web funcional (React
más Tailwind CSS), navegable de principio a fin: permite crear, listar, consultar y editar información real
en cada uno de los tres módulos, y simula de forma interactiva el comportamiento de guardado local y
sincronización que exige el enunciado.

\safefig{01-inicio.png}{Pantalla de inicio, con tres accesos directos a los módulos principales.}

La pantalla de inicio evita sobrecargar al usuario con paneles secundarios: muestra únicamente los tres
accesos a los módulos, con un conteo breve de contexto en cada uno, y un aviso de sincronización pendiente
cuando corresponde. Esto reduce la carga cognitiva del primer contacto (heurística de diseño estético y
minimalista): la jueza decide de inmediato qué tarea realizar, en vez de interpretar listas antes de poder
actuar.

% ---------------- Módulo 1 ----------------
\FloatBarrier
\section{Módulo 1: Registro de casos judiciales}

La jueza necesita registrar un nuevo caso, consultar los casos existentes y actualizar su avance. El diseño
facilita identificar rápidamente el estado de cada caso y reconocer cuáles requieren atención.

{\small
\begin{longtable}{@{}p{4.4cm}p{2.6cm}p{8.2cm}@{}}
\caption{Campos del módulo de casos judiciales.}\\
\toprule
\textbf{Campo} & \textbf{¿Obligatorio?} & \textbf{Validación / criterio} \\
\midrule
\endfirsthead
\multicolumn{3}{l}{\textit{(continuación)}} \\
\toprule
\textbf{Campo} & \textbf{¿Obligatorio?} & \textbf{Validación / criterio} \\
\midrule
\endhead
\bottomrule
\endfoot
Código del caso & --- & Autogenerado por la aplicación, de solo lectura. \\
Tipo de conflicto o materia & Obligatorio & Selección de lista controlada. \\
Descripción del motivo & Obligatorio & Mínimo 10 caracteres. \\
Estado & Obligatorio & En trámite / Concluido / Suspendido. Por defecto \enquote{En trámite}. \\
Observaciones & Opcional & --- \\
Partes: nombres y apellidos & Obligatorio & Al menos una parte. \\
Partes: documento de identidad & Opcional & \enquote{Si corresponde}, según el propio caso. \\
Partes: teléfono & Opcional & \enquote{Si se dispone}, puede no existir en zona rural. \\
Partes: dirección/comunidad & Obligatorio & Acepta solo el nombre de la comunidad. \\
Partes: rol en el caso & Obligatorio & Solicitante / Invitado / Testigo / Otro. \\
Seguimiento: fecha y descripción & Obligatorio al agregar & No se exige tener seguimiento para guardar el caso. \\
Seguimiento: resultado o acuerdo & Condicional & Obligatorio solo si el estado pasa a \enquote{Concluido}. \\
\end{longtable}
}

\safefig{02-casos-lista.png}{Listado de casos. El estado y el estado de sincronización se muestran con color, ícono y texto a la vez, nunca solo con color, por accesibilidad.}

\safefig{03-caso-nuevo.png}{Formulario de nuevo caso. Los campos obligatorios llevan asterisco; los opcionales, la etiqueta \enquote{(opcional)} explícita, de modo que un dato que no existe en campo rural nunca bloquea el registro.}

\safefig{04-caso-detalle.png}{Detalle de un caso, con la línea de tiempo de seguimiento. Marcar un caso como concluido exige registrar el resultado o acuerdo alcanzado y pide confirmación antes de aplicarse, evitando que una acción poco frecuente y difícil de revertir ocurra por error.}

% ---------------- Módulo 2 ----------------
\FloatBarrier
\section{Módulo 2: Registro de actuaciones notariales}

La jueza necesita dejar constancia de una actuación notarial y localizarla después. El formulario explica
con claridad qué información se requiere, evitando términos jurídicos o técnicos innecesarios.

{\small
\begin{longtable}{@{}p{4.4cm}p{2.6cm}p{8.2cm}@{}}
\caption{Campos del módulo de actuaciones notariales.}\\
\toprule
\textbf{Campo} & \textbf{¿Obligatorio?} & \textbf{Validación / criterio} \\
\midrule
\endfirsthead
\multicolumn{3}{l}{\textit{(continuación)}} \\
\toprule
\textbf{Campo} & \textbf{¿Obligatorio?} & \textbf{Validación / criterio} \\
\midrule
\endhead
\bottomrule
\endfoot
Código de registro & --- & Autogenerado por la aplicación, de solo lectura. \\
Tipo de actuación notarial & Obligatorio & Selección con lenguaje simple, sin jerga innecesaria. \\
Descripción o asunto & Obligatorio & Mínimo 10 caracteres. \\
Estado de la atención & Obligatorio & Pendiente / Atendida / Concluida. Por defecto \enquote{Pendiente}. \\
Fecha de atención & Opcional & Puede definirse después de la solicitud. \\
Personas: nombre & Obligatorio & Al menos una persona. \\
Personas: documento / teléfono & Opcional & \enquote{Si corresponde} / \enquote{si se dispone}. \\
Personas: dirección/comunidad & Obligatorio & Mismo criterio que en Casos. \\
Personas: participación & Obligatorio & Solicitante / Declarante / Testigo. \\
Resultado de la actuación & Condicional & Obligatorio solo si el estado es \enquote{Concluida}. \\
Documento adjunto o referencia & Opcional & Acepta referencia a un archivo físico. \\
\end{longtable}
}

\safefig{05-notarial-lista.png}{Listado de actuaciones notariales. Reutiliza la misma estructura visual y de navegación que el módulo de casos (heurística de consistencia y estándares), de modo que lo aprendido en un módulo se transfiere directamente al otro.}

El resultado de la actuación solo se exige cuando el estado pasa a \enquote{Concluida}; el resto del tiempo
el campo queda disponible pero no obligatorio, evitando registros a medio completar por exigir un dato que
aún no existe.

% ---------------- Módulo 3 ----------------
\FloatBarrier
\section{Módulo 3: Agenda o calendario}

La jueza necesita organizar audiencias, reuniones, visitas y otras actividades, y reconocer con facilidad si
tiene tareas próximas.

{\small
\begin{longtable}{@{}p{4.4cm}p{2.6cm}p{8.2cm}@{}}
\caption{Campos del módulo de agenda.}\\
\toprule
\textbf{Campo} & \textbf{¿Obligatorio?} & \textbf{Validación / criterio} \\
\midrule
\endfirsthead
\multicolumn{3}{l}{\textit{(continuación)}} \\
\toprule
\textbf{Campo} & \textbf{¿Obligatorio?} & \textbf{Validación / criterio} \\
\midrule
\endhead
\bottomrule
\endfoot
Título de la actividad & Obligatorio & --- \\
Fecha & Obligatorio & --- \\
Hora de inicio / término & Opcional & Permite actividades sin horario fijo. \\
Lugar o comunidad & Obligatorio & Necesario para planificar el desplazamiento. \\
Descripción / motivo & Opcional & --- \\
Estado & Obligatorio & Programada / Realizada / Cancelada. \\
Categoría & Obligatorio & Define color e ícono; nunca solo color. \\
\end{longtable}
}

\safefig{06-agenda-mes.png}{Vista mensual de la agenda. El día actual queda resaltado y cada actividad muestra una etiqueta de color por categoría siempre acompañada de texto. Se agregó un buscador de fecha específica junto a la navegación por mes, para no depender de desplazarse mes a mes cuando se conoce la fecha exacta buscada.}

% ---------------- Sin conexión ----------------
\FloatBarrier
\section{Funcionamiento sin conexión}

\safefig{07-sin-conexion.png}{Barra superior en estado \enquote{Sin conexión}, con el registro trabajando exclusivamente en el dispositivo.}

Todo registro nuevo o editado se guarda de inmediato en el dispositivo. Si hay conexión, queda marcado como
\enquote{Sincronizado}; si no la hay, se marca como \enquote{Guardado en el dispositivo} (pendiente), sin
bloquear ni interrumpir el trabajo de la jueza. Apenas la aplicación detecta que la conexión volvió,
sincroniza los pendientes de forma automática, sin exigir que la jueza recuerde hacerlo manualmente ---en la
práctica, ella no necesariamente sabrá en qué momento exacto la tableta recuperó señal. El botón
\enquote{Sincronizar ahora} se mantiene visible como control manual de respaldo, y la hora de la última
sincronización queda siempre a la vista en la barra superior, para que la confianza en el estado de los
datos no dependa de la memoria del usuario. Al sincronizar, la aplicación confirma cuántos registros se
sincronizaron correctamente, tal como lo exige el enunciado del caso.

% ---------------- Confirmación de cambios ----------------
\FloatBarrier
\section{Confirmación de cambios al editar}

\safefig{08-confirmar-cambios.png}{Diálogo mostrado al guardar la edición de un registro existente, con el valor anterior y el nuevo de cada campo modificado.}

Cuando se permite editar información, el enunciado exige que la interfaz indique qué datos se están
modificando y confirme que los cambios quedaron guardados. Por eso, al editar cualquier registro existente,
antes de aplicar el cambio la aplicación muestra un resumen con cada campo modificado (valor anterior y
valor nuevo) y exige una confirmación explícita (\enquote{Confirmar y guardar}) o permite volver a editar
(\enquote{Seguir editando}). Si no se detecta ningún cambio real, la aplicación lo indica y no interrumpe
con un diálogo innecesario.

% ---------------- Principios y heurísticas ----------------
\FloatBarrier
\section{Principios de diseño, heurísticas y factores humanos aplicados}

Cada decisión de diseño relevante del prototipo se documenta también dentro de la propia aplicación
(pantalla \enquote{Justificación de diseño}), relacionándola con un principio de Norman, una heurística de
usabilidad de Nielsen o un factor humano concreto, y explicando qué necesidad del usuario atiende:

{\small
\begin{longtable}{@{}p{5cm}p{4.3cm}p{6.2cm}@{}}
\caption{Principios, heurísticas y factores humanos aplicados en el prototipo.}\\
\toprule
\textbf{Decisión de diseño} & \textbf{Principio / heurística / factor humano} & \textbf{Necesidad que atiende} \\
\midrule
\endfirsthead
\multicolumn{3}{l}{\textit{(continuación)}} \\
\toprule
\textbf{Decisión de diseño} & \textbf{Principio / heurística / factor humano} & \textbf{Necesidad que atiende} \\
\midrule
\endhead
\bottomrule
\endfoot
Inicio con 3 accesos directos grandes a los módulos, sin listas secundarias & Heurística 8 --- Diseño estético y minimalista & El usuario decide de inmediato qué tarea hacer, sin interpretar listas antes de poder actuar. \\
Barra de estado de conexión y sincronización siempre visible & Heurística 1 --- Visibilidad del estado del sistema & Saber en todo momento si el trabajo está seguro. \\
Código de caso/actuación autogenerado, de solo lectura & Prevención de error humano (Norman) & Evita duplicados y errores de tipeo en identificadores. \\
Lenguaje simple en formularios, sin jerga jurídica innecesaria & Heurística 2 --- Coincidencia entre el sistema y el mundo real & Usuario sin experiencia frecuente en apps digitales. \\
Asterisco + etiqueta \enquote{(opcional)} explícita en cada campo & Heurística 5 --- Prevención de errores + Visibilidad (Norman) & Evita bloqueos por datos que no existen en zona rural. \\
Validación en línea con mensaje específico antes de enviar & Heurística 9 --- Ayudar a reconocer y recuperarse de errores & Reduce frustración frente a intentos repetidos a ciegas. \\
Confirmación visible (\enquote{toast}) tras guardar cada registro & Feedback (Norman) & Cumple lo pedido explícitamente por el caso. \\
Estado con color + ícono + texto, nunca solo color & Accesibilidad (daltonismo) + reconocimiento vs.\ recuerdo & Reconocer el estado de un caso de un vistazo. \\
Ícono siempre acompañado de texto en toda acción & Heurística 6 --- Reconocimiento en vez de recuerdo & El usuario no memoriza símbolos abstractos. \\
Botón \enquote{+ Nuevo} grande y fijo por módulo & Affordance clara + Mapping directo & La acción más frecuente siempre alcanzable con el pulgar. \\
Confirmación con opción de cancelar antes de concluir un caso & Heurística 3 --- Control y libertad del usuario & Previene una acción difícil de revertir por error. \\
Al editar, diálogo que lista los campos modificados (antes $\rightarrow$ después) antes de guardar & Visibilidad del estado + Heurística 3 --- Control y libertad del usuario & Cumple lo pedido explícitamente: indicar qué datos se están modificando, no solo confirmar al final. \\
Misma estructura de navegación y formularios en los 3 módulos & Heurística 4 --- Consistencia y estándares & El modelo mental de un módulo sirve para los otros dos. \\
Estados vacíos con instrucción, nunca pantalla en blanco & Guía activa ante la incertidumbre & Reduce la ansiedad del primer uso. \\
Tipografía grande, alto contraste, objetivos táctiles amplios & Factor humano: uso táctil, exteriores con sol & Legibilidad y precisión de toque en campo. \\
Navegación lateral fija pensada para tableta en horizontal & Factor humano: tamaño y orientación real del dispositivo & Aprovecha el espacio de 10 pulgadas sin saturarlo. \\
\enquote{Sincronizar ahora} con resultado explícito de éxito & Visibilidad del estado + literal del caso & Confirma si el proceso se completó correctamente. \\
Badge ámbar \enquote{Guardado en el dispositivo} vs.\ verde \enquote{Sincronizado} & Visibilidad + reduce carga de memoria & El usuario nunca tiene que adivinar si algo se sincronizó. \\
Sincronización automática al recuperar conexión + botón manual de respaldo & Automatización sin perder control (Heurística 3) & No depende de que la jueza recuerde sincronizar manualmente. \\
Hora de la última sincronización siempre visible en la barra superior & Heurística 1 --- Visibilidad del estado del sistema & Confianza en que la información está al día, sin tener que preguntarse cuándo se guardó. \\
\end{longtable}
}

% ---------------- Limitaciones ----------------
\FloatBarrier
\section{Limitaciones del prototipo}

\begin{itemize}[leftmargin=1.3em]
  \item No hay autenticación real ni backend: la sincronización es simulada.
  \item No se implementa carga real de archivos; los documentos quedan como referencia de nombre.
  \item El acceso desde computadora reutiliza el mismo diseño responsivo, sin un layout de escritorio dedicado.
  \item No se realizaron pruebas de usabilidad con usuarios reales.
\end{itemize}

% ---------------- Conclusión ----------------
\FloatBarrier
\section{Conclusión}

El prototipo cubre los tres módulos pedidos por el caso, permite trabajar sin conexión indicando con
claridad el estado de cada registro, y documenta explícitamente el principio o heurística detrás de cada
decisión de diseño. Las decisiones priorizan de forma sistemática la condición real de uso descrita en el
enunciado: una persona con baja familiaridad con aplicaciones digitales, operando desde una tableta táctil,
en un entorno donde la conexión a Internet es la excepción y no la norma.

\end{document}
